% In this file you should put the actual content of the blueprint.
% It will be used both by the web and the print version.
% It should *not* include the \begin{document}
%
% If you want to split the blueprint content into several files then
% the current file can be a simple sequence of \input. Otherwise It
% can start with a \section or \chapter for instance.

\section{Overview}

This blueprint documents the Lean~4 formalization of the ACTUS technical
specification.  It covers all eighteen contract types of the standard
(\S7.1--7.18).  The stateful lending family (PAM, LAM, NAM, ANN, CLM) is modeled
twice --- relationally (the inductive \texttt{Step}, the readable spec) and
functionally (computable state-transition and payoff functions), with the two
proven to agree.  The remaining types --- positions, derivatives, swaps, credit
enhancement, and the exotic amortizer --- are direct cash-flow builders over the
shared schedule-generation engine.

The whole development is generic over the amount type and instantiated twice:
at \texttt{Float} for the executable engine --- validated against the ACTUS
reference test suite, where it reproduces all $236$ gated reference contracts
($3155$ cash flows) exactly --- and at the real numbers $\mathbb{R}$ for the
quantitative metatheorems below.  The order/field structure those bounds need
exists over $\mathbb{R}$ but not over \texttt{Float} (where $\mathtt{NaN}$ even
breaks $a \le a$), so they are genuine theorems over $\mathbb{R}$ rather than
empirical checks.

\section{Utility layer}

\begin{definition}[Contract Role Sign $R$]\label{def:sign}
        \lean{Actus.Util.Conventions.sign}
        \leanok
        The contract-role sign $R : \mathtt{CNTRL} \to \{-1, +1\}$ (\S3.7),
        mapping a role to $+1$ for a claim (asset) or $-1$ for an obligation.
\end{definition}

\begin{definition}[Year Fraction $Y$]\label{def:yearFraction}
        \lean{Actus.Util.DayCount.yearFraction}
        \leanok
        The year-fraction convention $Y : s, t, \mathtt{DCC} \to \mathbb{R}$
        (\S3.6), giving the fraction of a year between two dates under a
        day-count convention (Actual/360, Actual/365, 30E/360, Actual/Actual
        ISDA, Business/252).
\end{definition}

\begin{definition}[Schedule function $S$]\label{def:schedule}
        \lean{Actus.Util.Schedule.schedule}
        \uses{def:sign, def:yearFraction}
        \leanok
        The schedule function $S(s, c, T, B)$ (\S3.1): the cyclic event times
        from anchor $s$ with cycle $c$ up to end date $T$, with end-of-month and
        business-day conventions applied.
\end{definition}

\begin{definition}[Annuity Amount $A$]\label{def:annuity}
        \lean{Actus.Util.Schedule.annuity}
        \uses{def:yearFraction}
        \leanok
        The annuity amount $A(s,T,n,a,r)$ (\S3.8) sizing the constant ANN
        instalment so the notional is fully repaid by maturity.
\end{definition}

\section{Contract types}

\begin{definition}[PAM state transition]\label{def:pam-stf}
        \lean{Actus.Contract.PAM.stf}
        \uses{def:sign, def:yearFraction}
        \leanok
        The Principal-at-Maturity state-transition function (\S7.1): one branch
        per event type (IED, IP, MD, RR, \dots).
\end{definition}

\begin{definition}[PAM relational step]\label{def:pam-step}
        \lean{Actus.Contract.PAM.Step}
        \uses{def:pam-stf}
        \leanok
        The relational one-step transition; each constructor's target is the
        functional next-state \texttt{stf}.
\end{definition}

\begin{definition}[LAM / NAM / ANN]\label{def:amortizers}
        \lean{Actus.Contract.LAM.stf}
        \uses{def:pam-stf, def:annuity}
        \leanok
        The amortizing contracts, defined as deltas over PAM: a linear
        principal redemption (LAM), interest-first redemption (NAM), and
        constant-instalment annuity with rate-reset re-amortization (ANN).
\end{definition}

\begin{definition}[CLM: Call Money]\label{def:clm}
        \lean{Actus.Contract.Execution.clmCashflows}
        \uses{def:pam-stf}
        \leanok
        Call money (\S7.6): PAM state transitions with interest capitalizing
        (\texttt{IPCI}) until maturity.  The open-maturity ``call'' has no fixed
        maturity --- it closes on an observed exercise ($\mathtt{XD} =
        \sup\tau(O^{ev}(\mathtt{CID}))$), settling the grown notional as
        \texttt{STD} after the exercise-notice period.
\end{definition}

\subsection{Schedule generation and execution}

\begin{definition}[Event schedule generation]\label{def:gensched}
        \lean{Actus.Contract.Execution.genSchedule}
        \uses{def:schedule}
        \leanok
        \texttt{genSchedule} builds the full event \texttt{Schedule} (IED,
        IP/IPCI, PR, RR/RRF, FP, SC, MD, \dots) for a stateful contract from its
        dictionary terms, following the Contract-Schedule tables (\S7.1--7.6);
        \texttt{runSchedule} then folds the contract's \texttt{stf}/\texttt{pof}
        over it.  This is the type-agnostic core shared by every contract family.
\end{definition}

\subsection{Stateless cash-flow builders}

\begin{definition}[Position contracts]\label{def:position}
        \lean{Actus.Contract.Execution.comCashflows,
              Actus.Contract.Execution.stkCashflows}
        \uses{def:sign}
        \leanok
        COM (commodity, \S7.10) and STK (stock, \S7.9): a purchase (\texttt{PRD})
        and termination (\texttt{TD}) at the observed price, plus observed
        dividends (\texttt{DV}) for stock.  (CSH, \S7.8, has no scheduled cash
        flow.)
\end{definition}

\begin{definition}[Derivatives on an underlying]\label{def:derivative}
        \lean{Actus.Contract.Execution.optnsCashflows,
              Actus.Contract.Execution.futurCashflows,
              Actus.Contract.Execution.fxoutCashflows,
              Actus.Contract.Execution.capflCashflows}
        \uses{def:sign, def:yearFraction}
        \leanok
        Cash-settled derivatives: a European option's intrinsic value (OPTNS,
        \S7.15), a future's linear payoff $S - F$ (FUTUR, \S7.16), an FX
        outright's two-currency settlement (FXOUT, \S7.11), and an interest-rate
        cap/floor's rate-excess over a child contract's reset rate (CAPFL,
        \S7.14).
\end{definition}

\begin{definition}[Swaps]\label{def:swap}
        \lean{Actus.Contract.Execution.swppvCashflows,
              Actus.Contract.Execution.swapsCashflows}
        \uses{def:amortizers}
        \leanok
        SWPPV (plain-vanilla interest-rate swap, \S7.12) pays a fixed and a
        floating interest leg per period; SWAPS (\S7.13) is a composite of two
        child legs run by the ordinary engine and combined gross or net.
\end{definition}

\begin{definition}[Credit enhancement]\label{def:credit}
        \lean{Actus.Contract.Execution.cegCashflows,
              Actus.Contract.Execution.cecCashflows}
        \uses{def:sign}
        \leanok
        CEG (guarantee, \S7.17) and CEC (collateral, \S7.18) over covered
        contracts: a premium at purchase, and on a covered contract's observed
        credit event the covered exposure ($\sum$ notional, plus accrued interest
        when guaranteed) settles --- capped at the observed collateral value for
        CEC.
\end{definition}

\begin{definition}[Exotic amortizer and deposit]\label{def:laxump}
        \lean{Actus.Contract.Execution.laxCashflows,
              Actus.Contract.Execution.umpCashflows}
        \uses{def:amortizers, def:annuity}
        \leanok
        LAX (exotic amortizer, \S7.3) redeems (\texttt{PR}) or draws (\texttt{PI})
        principal on piecewise \emph{array} schedules; UMP (undefined-maturity
        profile, \S7.7) is a non-maturity deposit whose interest capitalizes,
        settling at the agreed price plus accrued interest on termination.
\end{definition}

\section{Metatheorems}

\begin{theorem}[Relational/functional agreement]\label{thm:agree}
        \lean{Actus.Contract.Agree.pam_step_to_fun}
        \uses{def:pam-step}
        \leanok
        Every relational PAM step lands on the functional next-state for some
        event and time; conversely every admissible event/time yields a step.
        The relation is exactly the graph of \texttt{stf}.
\end{theorem}

\begin{theorem}[Status-date monotonicity]\label{thm:mono}
        \lean{Actus.Contract.Properties.pam_trace_mono}
        \uses{thm:agree}
        \leanok
        Along any execution trace the status date is non-decreasing: each step
        requires $\mathtt{Sd} \le t$ and sets $\mathtt{Sd} := t$.
\end{theorem}

\begin{theorem}[Maturity]\label{thm:maturity}
        \lean{Actus.Contract.Properties.pam_md_nt}
        \uses{def:pam-stf}
        \leanok
        The maturity event zeroes the notional principal (and accrued interest
        and fees); termination additionally zeroes the interest rate.
\end{theorem}

\subsection{Quantitative bounds and conservation (over $\mathbb{R}$)}

\begin{theorem}[Rate cap/floor bound]\label{thm:rate-bound}
        \lean{Actus.Contract.Properties.pam_rr_rate_mem,
              Actus.Contract.Properties.ann_rr_rate_mem,
              Actus.Contract.Properties.clamp_mem}
        \uses{def:pam-stf, def:amortizers}
        \leanok
        After a rate reset the nominal rate lies within the life floor/cap
        window $[\mathtt{RRLF}, \mathtt{RRLC}]$ (whenever both are present and
        well-ordered).  Holds for PAM and, by reuse of PAM's rate logic, for
        ANN.  With every cap/floor absent the reset instead returns the raw
        repricing target $O^{rf}(\mathtt{RRMO},t)\cdot\mathtt{RRMLT}+\mathtt{RRSP}$
        (\texttt{pam\_rr\_uncapped}).
\end{theorem}

\begin{theorem}[Redemption non-overshoot]\label{thm:no-overshoot}
        \lean{Actus.Contract.Properties.redeemed_no_overshoot,
              Actus.Contract.Properties.lam_pr_nt_mem,
              Actus.Contract.Properties.redeemed_abs_le}
        \uses{def:amortizers}
        \leanok
        A principal redemption never overshoots: with a sign-normalised
        nonnegative state the post-redemption notional stays in $[0, \mathtt{Nt}]$,
        and the redeemed amount never exceeds the outstanding notional in
        magnitude ($|\,\mathrm{redeemed}\,| \le |\mathtt{Nt}|$).
\end{theorem}

\begin{theorem}[Strict amortization and full repayment]\label{thm:amortization}
        \lean{Actus.Contract.Properties.lam_pr_nt_lt,
              Actus.Contract.Properties.lam_pr_full_repay}
        \uses{thm:no-overshoot}
        \leanok
        With a strictly positive notional and instalment a redemption strictly
        decreases the notional (guaranteed progress); once the instalment
        reaches the outstanding notional the step drives it to exactly $0$
        (clean repayment).
\end{theorem}

\begin{theorem}[Interest capitalization conserves value]\label{thm:ipci}
        \lean{Actus.Contract.Properties.pam_ipci_capitalizes}
        \uses{def:pam-stf}
        \leanok
        The \texttt{IPCI} event moves accrued interest into the notional and
        resets accrued interest to $0$: value is relocated, not created.  (This
        one holds definitionally at any amount type.)
\end{theorem}

\begin{definition}[Principal-redemption trace]\label{def:prtrace}
        \lean{Actus.Contract.Properties.PRTrace}
        \uses{def:amortizers}
        \leanok
        \texttt{PRTrace} is the sub-relation of the LAM execution trace
        consisting of zero or more admissible principal-redemption (\texttt{PR})
        steps --- the regime in which the notional should be conserved (events
        like \texttt{IED}/\texttt{IPCI} reset or grow it).
\end{definition}

\begin{theorem}[Trace-level conservation of the notional]\label{thm:conservation}
        \lean{Actus.Contract.Properties.prTrace_nt_mem,
              Actus.Contract.Properties.prTrace_abs_le,
              Actus.Contract.Properties.prTrace_isTrace}
        \uses{def:prtrace, thm:no-overshoot}
        \leanok
        Along an entire principal-redemption trace the notional stays in
        $[0, \mathtt{Nt}_0]$ and never increases in magnitude --- the
        single-step non-overshoot lifted to whole executions by induction over
        the trace.  Every \texttt{PRTrace} is moreover a genuine LAM execution
        trace, so this is a statement about real runs.
\end{theorem}

\subsection{Payoff bounds for the derivative and credit families (over $\mathbb{R}$)}

\begin{theorem}[Cap/floor and option intrinsics are non-negative]\label{thm:intrinsic}
        \lean{Actus.Contract.Properties.capfl_intrinsic_nonneg,
              Actus.Contract.Properties.option_call_intrinsic_nonneg,
              Actus.Contract.Properties.option_put_intrinsic_nonneg,
              Actus.Contract.Properties.option_call_pos_iff}
        \uses{def:derivative}
        \leanok
        The cap/floor rate-excess kernel
        $\max(r - \mathtt{cap}, 0) + \max(\mathtt{floor} - r, 0)$ and the
        European call/put intrinsics $\max(S - K, 0)$, $\max(K - S, 0)$ are all
        $\ge 0$, so a long cap, floor, or option never settles negative; a call
        is in the money exactly when the underlying exceeds the strike.
\end{theorem}

\begin{theorem}[Settlement conservation for FX and swaps]\label{thm:conservation-deriv}
        \lean{Actus.Contract.Properties.fxout_delivery_sum,
              Actus.Contract.Properties.swppv_net_eq_legs}
        \uses{def:swap, def:derivative}
        \leanok
        An FX outright's two delivery legs net to the par difference
        $R \cdot (\mathtt{Nt_1} - \mathtt{Nt_2})$ at unit FX; a swap's per-period
        fixed and floating legs net to $R \cdot N \cdot (\mathtt{fix} -
        \mathtt{flt}) \cdot Y$ --- exactly what the net-settlement fold computes
        from the two gross legs.
\end{theorem}

\begin{theorem}[Credit-enhancement payout bounds]\label{thm:credit-bound}
        \lean{Actus.Contract.Properties.cec_payout_le_value,
              Actus.Contract.Properties.cec_payout_le_claim,
              Actus.Contract.Properties.ceg_exposure_nonneg,
              Actus.Contract.Properties.foldl_add_nonneg}
        \uses{def:credit}
        \leanok
        A collateral payout $\min(\mathtt{coverage}\cdot\mathtt{exposure},
        \mathtt{value})$ never exceeds either the posted collateral value or the
        covered claim; the covered exposure, a sum of per-leg magnitudes, is
        non-negative, so the payout carries the contract-role sign.
\end{theorem}
